\documentclass[12pt,letterpaper]{article}
\usepackage[margin=1in]{geometry}
\usepackage{amsmath,amssymb,amsthm,blkarray,braket,caption,enumitem,fancyvrb,graphicx,hyperref,relsize,verbatim}
\usepackage[numbers]{natbib}
\usepackage[section]{placeins}
\graphicspath{{./images/}}
\bibliographystyle{apalike}

\begin{document}

\title{Chapter 1 - Introduction to Groups}
\author{Emil Marinov}
\maketitle

\section{Problem 1}
\textit{Determine which of the following binary operations are associative:}

\begin{enumerate}
   \item \textit{The operation $*$ on $\mathbb{Z}$ defined by $a * b = a - b$.} \\
   \\

   Proof by contradiction.
   Assume the operation $*$ is associative.
   Then, $(a * b) * c = a * (b * c) \forall a, b, c \in \mathbb{Z}$.

   \begin{equation}
      \begin{aligned}
         (a * b) * c = (a - b) - c = a - b - c \\
         a * (b * c) = a - (b - c) = a - b + c
      \end{aligned}
   \end{equation}

   This is a contradiction, as for any fixed $a$ and $b$, $\exists c : a - b - c \neq a - b + c$.
   Take any $c \neq 0$, for example.
   Thus, the operation $*$ is \textbf{not associative}.

   \item \textit{The operation $*$ on $\mathbb{R}$ defined by $a * b = a + b + ab$.} \\
   \\
   For $a, b, c \in \mathbb{R}$,

   \begin{equation}
      \begin{aligned}
         (a * b) * c = (a + b + ab) * c \\
         = (a + b + ab) + c + (a + b + ab)c \\
         = a + b + ab + c + ac + bc + abc \\
         = a + b + c + bc + ab + ac + abc \\
         = a + (b + c + bc) + a(b + c + bc) \\
         = a * (b + c + bc) \\
         = a * (b * c)
      \end{aligned}
   \end{equation}

   $\forall a, b, c \in \mathbb{R}$, $(a * b) * c = a * (b * c)$.
   Thus, the operation $*$ is \textbf{associative}.

   \textbf{Version 2 - The Imposter Additive Operation} \\
   Notice that $a + b + ab = (a + 1)(b + 1) - 1$ and define the map $\varphi : (\mathbb{R}, *) \to (\mathbb{R}^\times, \circ)$ as 
   $\phi(a) = a + 1$, or a linear translation over $\mathbb{R}$.

   \begin{equation}
      \begin{aligned}
         \phi(a * b) = 1 + (a * b) \\
         = 1 + (a + b + ab) \\
         = (a + 1)(b+1) \\
         = \phi(a) \circ \phi(b)
      \end{aligned}
   \end{equation}

   So, $\varphi$ is a homomorphism.
   It must now be determined if $\varphi$ is bijective.
   Proof by contradiction is used to show injection.
   Assume $a, b \in \mathbb{R} : a \neq b$ map to the same element $s \in \mathbb{R}$ in the range.
   
   \begin{equation}
      \begin{aligned}
         \varphi(a) = \varphi(b) = s \\
         a + 1 = b + 1 = s \\
         a = b = s - 1
      \end{aligned}
   \end{equation}

   This shows $a = b$, reaching a contridiction, and indicating that the assumption of $a \neq b$ is false.
   Proof by construction is used to prove surjectivity.
   For all elements $c \in (\mathbb{R}, \circ)$, there exists an element that maps to it, namely a left translation by 1.
   Because $\varphi$ is bijective, it is an isomorphic map, and there exists an inverse
   $(\mathbb{R}^\times, \circ) \to (\mathbb{R}, *)$ defined as $\phi^{-1}(b) = b - 1$.

   \begin{equation}
      \begin{aligned}
         (a * b) * c = \varphi^{-1}\left[ ((\varphi(a * b)) \circ \varphi(c) ) \right] \\
         = \varphi^{-1}\left[ ((\varphi(a) \circ \varphi(b)) \circ \varphi(c) ) \right] \\
         = \varphi^{-1}\left[ (\varphi(a) \circ (\varphi(b) \circ \varphi(c)) ) \right] \\
         = \varphi^{-1}\left[ (\varphi(a) \circ \varphi(b * c) ) \right] \\
         = \varphi^{-1}\left[ (\varphi(a) * \varphi(b * c) ) \right] \\
         = \varphi^{-1}\left[ (\varphi(a) * \varphi(b * c) ) \right] \\
         a * (b * c)
      \end{aligned}
   \end{equation}

   \begin{equation}
      \begin{aligned}
         (a*b)*c = a*(b*c) \quad \forall a, b, c \in \mathbb{R} && \text{By associative property of multiplication.}
      \end{aligned}
   \end{equation}

\end{enumerate}

\section{Problem 5}
\textit{Prove for all $n > 1$ that $\mathbb{Z}/n\mathbb{Z}$ is not a group under multiplication of residue classes.} \\
\\
%The existence of a single zero divisor, $\bar{0}$, without a left inverse indicates it is not a group.
%As the composition of $\bar{0}$ with any other element of the set is not injective, there exists no inverse.
%For injectivity to hold, $f(x) = f(y) \imply x = y$.
%However, $f(\bar{0}) = f(\bar{1})$, and $0 \neq 1$, in all cases where $n > 1$.
%Thus, the operation is not bijective, and non-invertible.

The left and right cancellation laws must hold for a group $G$.

\begin{equation}
   \begin{aligned}
      au = av \implies u = v \quad \forall a,u,v \in G && \text{By left cancellation.} \\
      ua = va \implies u = v \quad \forall a,u,v \in G && \text{By right cancellation.}
   \end{aligned}
\end{equation}

However, for $\mathbb{Z}/n\mathbb{Z}$, $\bar{0} * \bar{a} = \bar{0} * \bar{b}$, even when $a \neq b$.
Thus, $\mathbb{Z}/n\mathbb{Z}$ for $n > 1$ is \textbf{not a group}.

\end{document}